\documentclass{article}
\usepackage{amsmath,amssymb,amsthm}
\usepackage{algorithm,algorithmic}
\usepackage{hyperref}
\usepackage{geometry}
\geometry{margin=1in}
\newtheorem{theorem}{Theorem}
\newtheorem{lemma}{Lemma}
\newtheorem{assumption}{Assumption}
\newtheorem{remark}{Remark}
\begin{document}

\title{Warm‑Restart Stochastic Gradient Descent (SGDR) \\ Detailed Algorithmic Description and Convergence Theory}
\author{}
\date{}
\maketitle

%%%%%%%%%%%%%%%%%%%%%%%%%%%%%
\section*{Algorithm Name}
\textbf{Stochastic Gradient Descent with Cosine Annealing and Warm Restarts (SGDR)}  

The method follows the classic SGD update but replaces the ordinary learning‑rate schedule by a cosine‑annealed schedule that is periodically reset (“warm restart”).

%%%%%%%%%%%%%%%%%%%%%%%%%%%%%
\section*{Mathematical Formulation}

Let $f:\mathbb{R}^d\rightarrow\mathbb{R}$ be the (possibly non‑convex) objective and let $g_t$ denote an unbiased stochastic gradient
\[
\mathbb{E}[\,g_t\mid w_t\,]=\nabla f(w_t),\qquad 
\mathbb{E}\big[\|g_t-\nabla f(w_t)\|^2\mid w_t\big]\le\sigma^2 .
\]

\paragraph{Learning‑rate schedule.}
For a fixed maximal learning rate $\eta_{\max}>0$, a restart period $M\in\mathbb{N}$ and an optional period‑multiplication factor $T_{\rm mul}\ge 1$, define
\[
\boxed{\;
\eta_t \;=\; \eta_{\max}\,
\frac{1}{2}\Bigl(1+\cos\Bigl(\pi\frac{t\bmod M}{M}\Bigr)\Bigr)
\;}
\tag{1}
\]
The cosine term decays from $\eta_{\max}$ to $0$ within $M$ iterations.  
After $M$ iterations a \emph{warm restart} occurs:
\[
\text{if } (t+1)\bmod M =0\;\Longrightarrow\;
\begin{cases}
\eta_{t+1}\gets\eta_{\max},\\[2mm]
M\gets\lceil T_{\rm mul}\,M\rceil .
\end{cases}
\tag{2}
\]

\paragraph{Parameter update.}
\[
\boxed{\;
w_{t+1}=w_t-\eta_t\,g_t,\qquad t=0,1,\dots
\;}
\tag{3}
\]

%%%%%%%%%%%%%%%%%%%%%%%%%%%%%
\section*{Pseudocode}
\begin{algorithm}[H]
\caption{SGDR}
\begin{algorithmic}[1]
\STATE \textbf{Input:} $w_0$, $\eta_{\max}>0$, initial period $M$, period multiplier $T_{\rm mul}\ge1$, max iterations $T$.
\STATE Set $t\leftarrow0$, $M_{\text{cur}}\leftarrow M$.
\WHILE{$t<T$}
    \STATE Compute stochastic gradient $g_t$ (e.g. on a minibatch).
    \STATE $\displaystyle \eta_t\leftarrow \eta_{\max}\frac{1}{2}\Bigl(1+\cos\bigl(\pi\frac{t\bmod M_{\text{cur}}}{M_{\text{cur}}}\bigr)\Bigr)$.
    \STATE $w_{t+1}\leftarrow w_t-\eta_t g_t$.
    \STATE $t\leftarrow t+1$.
    \IF{$t\bmod M_{\text{cur}}=0$}
        \STATE $M_{\text{cur}}\leftarrow\lceil T_{\rm mul}\,M_{\text{cur}}\rceil$ \COMMENT{warm restart}
    \ENDIF
\ENDWHILE
\STATE \textbf{Output:} final iterate $w_T$ (or averaged iterate).
\end{algorithmic}
\end{algorithm}

%%%%%%%%%%%%%%%%%%%%%%%%%%%%%
\section*{Dry Run Example}
Consider the one‑dimensional quadratic 
\[
f(w)=(w-3)^2,\qquad \nabla f(w)=2(w-3).
\]
Take $w_0=0$, $\eta_{\max}=0.1$, $M=3$ (no period multiplication), and use the \emph{exact} gradient $g_t=\nabla f(w_t)$.  

\begin{center}
\begin{tabular}{c|c|c|c}
$t$ & $\eta_t$ (by (1)) & $w_{t+1}=w_t-\eta_t g_t$ & $f(w_{t+1})$\\\hline
0 & $0.1\cdot\frac{1}{2}(1+\cos0)=0.1$ &
$0-0.1\cdot2(0-3)=0+0.6=0.6$ &
$(0.6-3)^2=5.76$\\[2mm]
1 & $0.1\cdot\frac{1}{2}\bigl(1+\cos\frac{\pi}{3}\bigr)
=0.1\cdot\frac{1}{2}(1+0.5)=0.075$ &
$0.6-0.075\cdot2(0.6-3)=0.6-0.075\cdot(-4.8)=0.6+0.36=0.96$ &
$(0.96-3)^2=4.1856$\\[2mm]
2 & $0.1\cdot\frac{1}{2}\bigl(1+\cos\frac{2\pi}{3}\bigr)
=0.1\cdot\frac{1}{2}(1-0.5)=0.025$ &
$0.96-0.025\cdot2(0.96-3)=0.96-0.025\cdot(-4.08)=0.96+0.102=1.062$ &
$(1.062-3)^2\approx3.764$\\
\end{tabular}
\end{center}
After three steps the learning rate would be reset to $0.1$ (warm restart) and the process repeats with a new period if $T_{\rm mul}>1$.

%%%%%%%%%%%%%%%%%%%%%%%%%%%%%
\section*{Assumptions}
\begin{assumption}[Smoothness]\label{ass:smooth}
$f$ is $L$‑smooth, i.e.
\[
\|\nabla f(x)-\nabla f(y)\|\le L\|x-y\|,\qquad\forall x,y\in\mathbb{R}^d .
\]
Equivalently,
\[
f(y)\le f(x)+\langle\nabla f(x),y-x\rangle+\frac{L}{2}\|y-x\|^2 .
\tag{4}
\]
\end{assumption}
\begin{assumption}[Lower Boundedness]\label{ass:lower}
There exists $f^\star>-\infty$ such that $f(w)\ge f^\star$ for all $w$.
\end{assumption}
\begin{assumption}[Unbiased stochastic gradients]\label{ass:unbiased}
\[
\mathbb{E}[g_t\mid w_t]=\nabla f(w_t),\qquad 
\mathbb{E}\big[\|g_t-\nabla f(w_t)\|^2\mid w_t\big]\le\sigma^2 .
\tag{5}
\]
\end{assumption}
\begin{assumption}[Bounded learning‑rate]\label{ass:eta}
The maximal learning rate satisfies
\[
0<\eta_{\max}\le\frac{1}{L}.
\tag{6}
\]
Consequently $\eta_t\le\eta_{\max}$ for all $t$.
\end{assumption}

%%%%%%%%%%%%%%%%%%%%%%%%%%%%%
\section*{Main Convergence Theorem}
\begin{theorem}[Convergence of SGDR]\label{thm:sgdr}
Let $\{w_t\}_{t\ge0}$ be generated by \eqref{3}–\eqref{2} under Assumptions \ref{ass:smooth}–\ref{ass:eta}.  
Define the average learning rate over $T$ iterations
\[
\bar\eta_T\;:=\;\frac{1}{T}\sum_{t=0}^{T-1}\eta_t .
\]
Then for any $T\ge1$,
\[
\frac{1}{T}\sum_{t=0}^{T-1}\mathbb{E}\big[\|\nabla f(w_t)\|^2\big]
\;\le\;
\underbrace{\frac{2\big(f(w_0)-f^\star\big)}{\bar\eta_T\,T}}_{\text{deterministic decay}}
\;+\;
\underbrace{L\,\eta_{\max}\,\sigma^2}_{\text{variance floor}} .
\tag{7}
\]
Consequently, if $\sigma=0$ (full‑gradient case) the right‑hand side decays as $O(1/T)$.  For stochastic gradients the bound converges to a neighbourhood of a stationary point whose radius is $O(L\eta_{\max}\sigma^2)$.
\end{theorem}

%%%%%%%%%%%%%%%%%%%%%%%%%%%%%
\section*{Theoretical Properties}
\begin{itemize}
\item \textbf{Stability.} Because $\eta_t\le\eta_{\max}\le1/L$, the descent inequality \eqref{4} guarantees that the objective never increases in expectation beyond the variance term.
\item \textbf{Hyper‑parameter sensitivity.}
    \begin{enumerate}
        \item Larger $\eta_{\max}$ speeds up the deterministic decay term $2(f(w_0)-f^\star)/(\bar\eta_T T)$ but also enlarges the variance floor $L\eta_{\max}\sigma^2$.
        \item The period $M$ influences $\bar\eta_T$: longer periods increase the proportion of small learning rates, reducing $\bar\eta_T$ and hence slowing convergence.
        \item Multiplying $M$ after each restart (common in practice) yields a non‑increasing $\bar\eta_T$, which empirically helps escape shallow local minima.
    \end{enumerate}
\item \textbf{Comparison to baselines.}
    \begin{itemize}
        \item With a \emph{fixed} learning rate $\eta$, the standard SGD bound is 
        $\frac{1}{T}\sum_{t}\mathbb{E}\|\nabla f(w_t)\|^2\le\frac{2(f(w_0)-f^\star)}{\eta T}+L\eta\sigma^2$.
        SGDR replaces the constant $\eta$ by the average $\bar\eta_T$, which is typically larger early on (faster progress) while still guaranteeing a finite variance floor.
        \item Compared with a pure \emph{decaying} schedule $\eta_t=\frac{c}{\sqrt{t}}$, SGDR enjoys a simple periodic structure and, under the same $\eta_{\max}$, yields the same $O(1/T)$ deterministic rate but with a smaller variance term because $\eta_t$ never exceeds $\eta_{\max}$.
    \end{itemize}
\end{itemize}

%%%%%%%%%%%%%%%%%%%%%%%%%%%%%
\section*{Discussion}
The theorem shows that SGDR inherits the classic $O(1/T)$ decay of SGD for smooth non‑convex problems, while the cosine annealing provides a practical way to allocate larger steps early and to “re‑energize’’ the optimizer after each restart.  The bound is tight up to constants; improving the variance term would require variance‑reduction techniques (e.g., SVRG) which can be combined with warm restarts.

%%%%%%%%%%%%%%%%%%%%%%%%%%%%%
\section*{Preliminaries (Definitions and Lemmas)}
\begin{lemma}[One‑step descent under $L$‑smoothness]\label{lem:descent}
For any $w_t$ and any learning rate $\eta_t$,
\[
\mathbb{E}\big[f(w_{t+1})\mid w_t\big]
\le f(w_t)-\eta_t\|\nabla f(w_t)\|^2
+\frac{L\eta_t^2}{2}\,\mathbb{E}\big[\|g_t\|^2\mid w_t\big].
\tag{8}
\]
\end{lemma}
\begin{proof}
Using $L$‑smoothness \eqref{4} with $y=w_{t+1}=w_t-\eta_t g_t$,
\[
f(w_{t+1})\le f(w_t)-\eta_t\langle\nabla f(w_t),g_t\rangle
+\frac{L\eta_t^2}{2}\|g_t\|^2 .
\]
Take conditional expectation and invoke unbiasedness \eqref{5}:
\[
\mathbb{E}[\langle\nabla f(w_t),g_t\rangle\mid w_t]
= \|\nabla f(w_t)\|^2 .
\]
Thus \eqref{8} follows. \hfill\qedsymbol
\end{proof}

\begin{lemma}[Bound on second moment of stochastic gradient]\label{lem:second}
Under Assumption \ref{ass:unbiased},
\[
\mathbb{E}\big[\|g_t\|^2\mid w_t\big]\le
\|\nabla f(w_t)\|^2+\sigma^2 .
\tag{9}
\]
\end{lemma}
\begin{proof}
Decompose $g_t=\nabla f(w_t)+(g_t-\nabla f(w_t))$, square, and use $\mathbb{E}[\,\cdot\,]=\|\nabla f(w_t)\|^2+\mathbb{E}\|g_t-\nabla f(w_t)\|^2$; the variance bound gives \eqref{9}. \hfill\qedsymbol
\end{proof}

%%%%%%%%%%%%%%%%%%%%%%%%%%%%%
\section*{Main Proof (Step‑by‑Step Derivation)}
We prove Theorem \ref{thm:sgdr}. Throughout the proof we annotate each non‑trivial manipulation with a step number.

\begin{proof}
\textbf{Step 1 (Apply Lemma \ref{lem:descent}).}  
From \eqref{8} and Lemma \ref{lem:second} we obtain, for every $t$,
\[
\begin{aligned}
\mathbb{E}[f(w_{t+1})\mid w_t]
&\le f(w_t)-\eta_t\|\nabla f(w_t)\|^2
+\frac{L\eta_t^2}{2}\big(\|\nabla f(w_t)\|^2+\sigma^2\big)\\[1mm]
&= f(w_t)-\underbrace{\Bigl(\eta_t-\frac{L\eta_t^2}{2}\Bigr)}_{\alpha_t}\,\|\nabla f(w_t)\|^2
+\frac{L\eta_t^2\sigma^2}{2}.
\end{aligned}
\tag{10}
\]
\textbf{Step 2 (Lower bound $\alpha_t$).}  
Because of Assumption \ref{ass:eta}, $\eta_t\le\eta_{\max}\le 1/L$, thus
\[
\eta_t-\frac{L\eta_t^2}{2}\;\ge\;\frac{\eta_t}{2}.
\tag{11}
\]
\textbf{Step 3 (Take total expectation).}  
Denote $\mathcal{F}_t$ the $\sigma$‑algebra generated by $\{w_0,\dots,w_t\}$.  Taking expectation of \eqref{10} and using \eqref{11} yields
\[
\mathbb{E}[f(w_{t+1})]\le 
\mathbb{E}[f(w_t)]-\frac{\eta_t}{2}\,\mathbb{E}\big[\|\nabla f(w_t)\|^2\big]
+\frac{L\eta_t^2\sigma^2}{2}.
\tag{12}
\]

\textbf{Step 4 (Telescoping sum).}  
Sum \eqref{12} for $t=0,\dots,T-1$:
\[
\mathbb{E}[f(w_T)]\le f(w_0)-\frac{1}{2}\sum_{t=0}^{T-1}\eta_t\,
\mathbb{E}\big[\|\nabla f(w_t)\|^2\big]
+\frac{L\sigma^2}{2}\sum_{t=0}^{T-1}\eta_t^2 .
\tag{13}
\]

\textbf{Step 5 (Use lower bound of $f$).}  
Since $f(w_T)\ge f^\star$ by Assumption \ref{ass:lower},
\[
f^\star\le f(w_0)-\frac{1}{2}\sum_{t=0}^{T-1}\eta_t\,
\mathbb{E}\big[\|\nabla f(w_t)\|^2\big]
+\frac{L\sigma^2}{2}\sum_{t=0}^{T-1}\eta_t^2 .
\tag{14}
\]

\textbf{Step 6 (Upper bound $\eta_t^2$ by $\eta_{\max}\eta_t$).}  
Because $0\le\eta_t\le\eta_{\max}$,
\[
\eta_t^2\le\eta_{\max}\eta_t\quad\forall t.
\tag{15}
\]
Insert \eqref{15} into \eqref{14}:
\[
f^\star\le f(w_0)-\frac{1}{2}\sum_{t=0}^{T-1}\eta_t\,
\mathbb{E}\big[\|\nabla f(w_t)\|^2\big]
+\frac{L\sigma^2\eta_{\max}}{2}\sum_{t=0}^{T-1}\eta_t .
\tag{16}
\]

\textbf{Step 7 (Isolate the gradient term).}  
Rearrange \eqref{16}:
\[
\sum_{t=0}^{T-1}\eta_t\,
\mathbb{E}\big[\|\nabla f(w_t)\|^2\big]
\le
2\big(f(w_0)-f^\star\big)
+L\eta_{\max}\sigma^2\sum_{t=0}^{T-1}\eta_t .
\tag{17}
\]

\textbf{Step 8 (Divide by $T\bar\eta_T$).}  
Recall $\bar\eta_T=\frac{1}{T}\sum_{t=0}^{T-1}\eta_t$.  Divide both sides of \eqref{17} by $T\bar\eta_T$:
\[
\frac{1}{T}\sum_{t=0}^{T-1}\mathbb{E}\big[\|\nabla f(w_t)\|^2\big]
\le
\frac{2\big(f(w_0)-f^\star\big)}{T\bar\eta_T}
+L\eta_{\max}\sigma^2 .
\tag{18}
\]

\textbf{Step 9 (Conclusion).}  
Equation \eqref{18} is precisely the bound \eqref{7} claimed in the theorem. \hfill\qedsymbol
\end{proof}

%%%%%%%%%%%%%%%%%%%%%%%%%%%%%
\section*{Rate Derivation (With Constants)}
The average learning rate $\bar\eta_T$ for the cosine schedule can be computed explicitly.  
Within a single period of length $M$,
\[
\frac{1}{M}\sum_{k=0}^{M-1}\eta_{\max}\frac{1}{2}\bigl(1+\cos(\pi k/M)\bigr)
=\eta_{\max}\frac{1}{2}\Bigl(1+\frac{2}{\pi M}\sin\frac{\pi}{2}\Bigr)
=\eta_{\max}\Bigl(\frac{1}{2}+\frac{1}{\pi M}\Bigr).
\]
If the algorithm runs for $T=Q M$ (integer number of periods) the average over $T$ iterations coincides with the per‑period average, i.e.
\[
\bar\eta_T=\eta_{\max}\Bigl(\frac{1}{2}+\frac{1}{\pi M}\Bigr).
\]
Plugging this into \eqref{7} yields the concrete rate
\[
\frac{1}{T}\sum_{t=0}^{T-1}\mathbb{E}\big[\|\nabla f(w_t)\|^2\big]
\le
\frac{4\big(f(w_0)-f^\star\big)}{\eta_{\max}\big(1+2/(\pi M)\big)T}
+L\eta_{\max}\sigma^2 .
\]
Hence the deterministic term decays as $O(1/T)$, while the stochastic term is independent of $T$.

%%%%%%%%%%%%%%%%%%%%%%%%%%%%%
\section*{Special Cases}
\begin{enumerate}
\item \textbf{Full‑gradient case ($\sigma=0$).}  
The bound reduces to
\[
\frac{1}{T}\sum_{t=0}^{T-1}\|\nabla f(w_t)\|^2
\le\frac{2\big(f(w_0)-f^\star\big)}{\bar\eta_T\,T},
\]
showing pure $O(1/T)$ convergence to a stationary point.

\item \textbf{Constant learning rate ($\eta_t\equiv\eta$).}  
Our analysis recovers the classical SGD bound
\[
\frac{1}{T}\sum_{t=0}^{T-1}\mathbb{E}\|\nabla f(w_t)\|^2
\le\frac{2\big(f(w_0)-f^\star\big)}{\eta T}+L\eta\sigma^2 .
\]

\item \textbf{Increasing period ($M_{k+1}=c\,M_k$).}  
If $M$ grows geometrically, the contribution of later periods to $\bar\eta_T$ diminishes, leading to a slightly larger deterministic term but the same variance floor $L\eta_{\max}\sigma^2$.
\end{enumerate}

%%%%%%%%%%%%%%%%%%%%%%%%%%%%%
\end{document}